\section{Conclusion}
\label{sec:conclusion}

We presented \kunai, a packet-filter DSL for the structures
pcap-filter cannot express, such as nested encapsulation, repeated
protocols, and variable-length lists and options. The compiler
resolves filter expressions against protocol definitions written in a
subset of P4, classifies each variable-length structure from the
shape of its parser when the definition is loaded, and lowers it to a
counted walk, a TLV walk, or a sentinel-terminated chain with a
compile-time bound and a bounds check before every packet read.
Because the lowering is derived from structure rather than protocol
names, a new protocol definition extends the filter language without
recompiling the tool. Across six kernels and both attach points, the
generated bytecode was accepted by the verifier, matched and rejected
packets as intended, and stayed within a few percent of pcap-filter
where both languages express a filter. \kunai complements existing
capture tools rather than replacing them. pcap-filter remains a good
fit for filters over fixed header fields, and \kunai takes over where
the target field moves from packet to packet.